\section{Annotation Rubric and Prompts}
\label{app:annotation}

\paragraph{Annotation schema.}
Each comment is tagged with nine fields. Six \emph{binary pragmatic}
fields:
\begin{itemize}\itemsep1pt
  \item \textbf{Implicit sentiment} ($1$ if the comment expresses
    sentiment without using sentiment-bearing lexical items, e.g.\
    via metaphor, rhetorical question, or comparison).
  \item \textbf{Sarcasm} ($1$ if the comment's surface polarity is
    the opposite of the intended polarity).
  \item \textbf{Irony} ($1$ if the comment expresses incongruity
    between expectation and reality, without necessarily flipping
    polarity).
  \item \textbf{Idiom / figurative} ($1$ if the comment uses a fixed
    Vietnamese idiomatic chunk or non-literal expression).
  \item \textbf{Code-switching} ($1$ if the comment contains at least
    one English content word other than universally adopted
    abbreviations such as ``OK'', ``Internet'').
  \item \textbf{Mocking} ($1$ if the comment ridicules a person,
    group, or behaviour; can co-occur with sarcasm).
\end{itemize}
Plus one \emph{intended polarity} field (positive / neutral /
negative) and one \emph{emotion} field on the UIT-VSMEC-7 scale.

\paragraph{Annotation prompt to annotators.}
Annotators received a 4-page rubric with $36$ illustrative comments
($6$ per pragmatic phenomenon) and a flowchart that asks them to
work outside-in: first decide intended polarity, then mark sarcasm
and irony if the intended polarity differs from the surface,
then mark mocking, then mark idiom and code-switching as orthogonal
form features. Annotators were instructed to leave a free-text
rationale if they marked sarcasm or mocking.

\paragraph{Rationale-generation prompt to GPT-4o-mini.}
\begin{quote}\itshape\small
You are a Vietnamese pragmatic-sentiment annotator. Given a
Vietnamese social-media comment and its gold labels (intended
polarity, six pragmatic flags), write a 1--2 sentence Vietnamese
explanation that names the lexical or contextual cues that justify
the labels. Do not repeat the comment verbatim. If the comment is
sarcastic, explicitly name the cue (e.g.\ trailing ``=))'',
exaggerated praise, lexical contrast with the discourse context).
Output only the explanation; do not output the labels.
\end{quote}

\paragraph{GPT-4o-mini classification prompt (8-shot).}
The prompt contains a fixed task description, eight in-context
demonstrations (one per common phenomenon plus two control
positives and one control negative), and the test comment. The
in-context demonstrations are drawn from the \method{} \emph{training}
set with adjudicated labels and are held fixed across all test
comments to make the baseline reproducible.

\paragraph{Inter-annotator agreement.}
Krippendorff's $\alpha$ by field:
implicit $0.72$, sarcasm $0.74$, irony $0.66$, idiom $0.69$,
code-switching $0.84$, mocking $0.61$.
Cohen's $\kappa$ on the polarity-after-sarcasm-flip decision is
$0.69$. Adjudication touched $14\%$ of comments overall, with the
adjudicator agreeing with the majority of the two annotators
$92\%$ of the time and overruling both annotators $1.3\%$ of the
time (rare but non-zero).

\paragraph{Annotator compensation and consent.}
Three annotators (two reviewers and one adjudicator) were L1
Vietnamese speakers paid \$15/hour for a $14$-week annotation pass.
Annotators provided written informed consent that included a
description of the dataset's intended research use and the option to
withdraw before publication; no annotator withdrew. The annotation
protocol was reviewed by the authors' institutional ethics office
before annotation began.
